% Generated draft: requires mathematical and source review.
\documentclass[11pt]{article}
\usepackage[margin=25mm]{geometry}
\usepackage{fontspec}
\setmainfont{Latin Modern Roman}
\usepackage{amsmath,amssymb,mathtools,unicode-math}
\setmathfont{Latin Modern Math}
\usepackage{enumitem,needspace,xurl,hyperref}
\setlength{\parindent}{0pt}
\setlength{\parskip}{5pt}
\setlength{\emergencystretch}{3em}
\newcommand{\N}{\mathbb{N}}
\newcommand{\Z}{\mathbb{Z}}
\newcommand{\Q}{\mathbb{Q}}
\newcommand{\R}{\mathbb{R}}
\newcommand{\C}{\mathbb{C}}
\newcommand{\F}{\mathbb{F}}
\newcommand{\A}{\mathbb{A}}
\newcommand{\PP}{\mathbb{P}}
\newcommand{\E}{\mathbb{E}}
\newcommand{\eps}{\varepsilon}
\newcommand{\dd}{\,\mathrm{d}}
\begin{document}
\section*{UnsolvedMath ID: 701663}
\textbf{Source problem code:} catalogue entry
\subsection{Definitions and mathematical statement}
Let \(G\) be a finite group. Its integral group ring is
\[
\mathbb ZG=\left\{\sum_{g\in G}a_g g:a_g\in\mathbb Z\right\},
\]
with multiplication induced by the group operation and distributivity. Define the augmentation homomorphism by
\[
\varepsilon:\mathbb ZG\longrightarrow\mathbb Z,
\qquad \varepsilon\!\left(\sum_{g\in G}a_g g\right)=\sum_{g\in G}a_g.
\]
The group of normalized units is
\[
\mathrm V(\mathbb ZG)=\{u\in(\mathbb ZG)^\times:\varepsilon(u)=1\}.
\]
For any group \(H\), define its finite-order spectrum by
\[
\operatorname{Spec}(H)=\{n\in\mathbb Z_{>0}:\text{there exists }h\in H\text{ with }\operatorname{ord}(h)=n\}.
\]
Here \(\operatorname{ord}(h)=n\) means that \(n\) is the least positive integer with \(h^n=1\). Infinite orders are excluded from this spectrum, consistently with the source's explicit formulation in terms of integer orders.

\textbf{Problem.} Is it true that, for every finite group \(G\),
\[
\operatorname{Spec}(G)=\operatorname{Spec}\bigl(\mathrm V(\mathbb ZG)\bigr)?
\]
Equivalently, for every finite group \(G\) and every \(n\in\mathbb Z_{>0}\), is
\[
\bigl(\exists u\in\mathrm V(\mathbb ZG):\operatorname{ord}(u)=n\bigr)
\quad\Longleftrightarrow\quad
\bigl(\exists g\in G:\operatorname{ord}(g)=n\bigr)?
\]
The source says ``for any integer \(n\)''; positive integers are used here because finite element orders are positive. The record describes this as a well-known problem and labels it open. It also reports an affirmative answer when \(G\) is solvable, attributed to M. Hertweck, \emph{Comm. Algebra} \textbf{36} (2008), 3585--3588. These are supplied source claims, not independently verified assessments. The proposer attribution is corrupted in the supplied record and cannot be reliably reconstructed.
\subsection{Short English statement}
Does every finite-order normalized unit of the integral group ring of a finite group have an order that also occurs among elements of the group, and conversely? The supplied record reports an affirmative answer for solvable groups and labels the general problem open; its proposer attribution is damaged.
\subsection{Sources}
\begin{itemize}
\item Local UnsolvedMath registry: supplied complete record with permanent ID \(701663\), upstream ID \(701663\), exact source code \texttt{catalogue entry}, and title \emph{Kourovka Notebook Problem 20.5}. Supplied upstream revision: \texttt{d371ee917ac9647fbb06cd6fc529deec6f3dda82}. The record was created and updated on 2026-10-16 and carries the status \texttt{open}.
\item The supplied background identifies \emph{The Kourovka Notebook (Unsolved Problems in Group Theory)}, 20th issue (2022), and the 21st issue edition (2026), \texttt{arXiv:1401.0300v48}; original problem number 20.5. The collection description names E. I. Khukhro and V. D. Mazurov as editors. These bibliographic details are supplied attribution and have not been independently checked.
\item The supplied statement attributes the affirmative answer for solvable groups to M. Hertweck, \emph{Comm. Algebra} \textbf{36} (2008), 3585--3588. No article title, DOI, or URL is supplied.
\item The proposer field is malformed, reading \texttt{A. B\{\textbackslash{}"a}. A displaced fragment under the background's discussion and literature heading reads \texttt{chle, L. Margolis\}}. The record gives proposed year 2022, but does not supply an uncorrupted proposer attribution.
\end{itemize}
\end{document}
